\section{Modularity and Extensibility}
\label{sec:eval-modularity}

The first research question asks whether sensing, analysis, decision, and intervention can be separated into modules with stable contracts, such that a new module connects without modifying the core.
We assess this at three levels: the static dependency structure, the locality of a change that adds a module, and the connection of external providers built against the contract.

At the assembly level the dependency direction is enforced by the build.
The core application project references only the domain project, and the infrastructure and host projects reference the core, so no reference points outward from the core toward infrastructure or transport.
A reference that violated this direction would be a compile error rather than a convention to be remembered.
We encode the rule as an executable test that inspects the referenced assemblies at runtime and asserts that the domain depends on no other project assembly and that the application core depends on no infrastructure, transport, or host assembly.
The test passes as part of the suite, which turns the modularity rule from an assertion in prose into a checked invariant.

The locality of a change measures how much of the system must be touched to extend it.
Adding an in-process intervention module requires one registration entry, because the runtime resolves modules from a registry by identifier and never names a concrete module, so no existing module and no part of the reading runtime is edited.
The out-of-process path is the more demanding test, and we exercise it with three external providers that run as separate processes and speak the documented module-provider protocol: a reference decision provider, a reference eye-movement analyser, and a connector that bridges the analysis pipeline of a concurrent, independent master's thesis at DTU to our \texttt{fixation-analysis} port.\footnote{\emph{Reading the Struggle: Using Machine Learning to Detect Reading Difficulty in Near Real-Time}, by Katarina Kraljević and Sree Keerthi Desu (DTU, in preparation). Their pipeline detects reader struggle (through cognitive load, reading style, and a RIPA2 score) from eye-tracking features; the connector that adapts it to our analysis contract was written by us and uses their code unmodified. As their thesis had not been submitted at the time of writing, it cannot yet be cited.}
This third provider is the strongest evidence available to us, because it integrates a real research pipeline that we neither wrote nor designed the contract around: its addition changed no file under the backend, the other team's code is used unmodified, and, most tellingly, the recorded evaluation sessions of this chapter were themselves analysed by it.
The fixation, saccade, and regression events reported in \cref{sec:eval-performance,sec:eval-context} were produced by that external provider rather than by the platform's built-in detector of \cref{sec:impl-sensing}, so an independent team's algorithm did not merely connect over the contract but served as the live analysis source for our own evaluation.
The malformed and unauthorised connection paths of the protocol are exercised by a handshake smoke test.

We report these as facts about the implemented system, and we read them against the modularity claim and prior work in \cref{ch:discussion}.
